% V. Методика на експерименталната проверка.
% ВАЖНО: този раздел описва ПЛАНА. Никакви числа тук. Числата живеят в
% Раздел VI и се вписват едва след като са измерени.
\section{Методика на експерименталната проверка}
\label{sec:method}

\subsection{Какво се проверява}

Проверката отговаря на три въпроса, поставени независимо един от друг. Първо,
дава ли системата правилен отговор там, където правилният отговор е известен
предварително. Второ, съвпадат ли класиранията на двата метода за
многокритериална оценка върху едни и същи данни, и ако не, къде се разминават.
Трето, колко устойчива е препоръката спрямо изменение на входните допускания.

Първият въпрос е за коректност и има еднозначен отговор. Вторият и третият са
измервания, чиито резултати не се знаят предварително и не се предполагат тук.

\subsection{Набор от сценарии}

Проверката се извършва върху три групи сценарии. Първата съдържа малки сценарии,
за които оптималният портфейл е изчислен ръчно чрез изчерпателно изброяване;
тяхната роля е да уловят грешка в модела, а не в решателя. Втората съдържа
сценарии, породени по зададено правило, с контролиран брой проекти и плътност на
графа на зависимостите; те служат за наблюдение на поведението при нарастващ
размер. Третата съдържа поне един сценарий, съставен по действителна практика на
планиране.

\TODO{опиши произхода на сценария от практиката, или, ако такъв не бъде получен,
запиши изрично, че третата група отпада, и не я замествай с измислен пример}

\subsection{Процедура}

За всеки сценарий се изпълнява следната процедура. Сценарият се зарежда и
инвариантите му се проверяват. Оценките се изчисляват по всеки от двата метода.
За всеки вектор от оценки задачата за подбор се решава веднъж с лакомата
евристика и веднъж с решателя. Резултатите се записват заедно с изразходваното
време. Анализът на чувствителността се изпълнява отделно, като бюджетът се
изменя в зададен интервал със зададена стъпка, а теглата на критериите се
изменят по същия начин.

Условие за валидност: резултатът на решателя не може да бъде по-лош от този на
евристиката. Нарушение на това условие се третира като грешка в модела и спира
експеримента, а не се отчита като резултат.

\TODO{фиксирай интервалите и стъпките за изменение на бюджета и теглата, както и
броя повторения за измерването на времето}

\subsection{Форма на резултатите}

Резултатите се представят в три вида. Таблица с по един ред на сценарий, която
съдържа броя проекти, броя избрани позиции, сумарната оценка по двата метода и
времето за решаване. Таблица на разминаванията между двата метода за
многокритериална оценка, с по един ред на всяко разминаване в класирането.
Графика на устойчивостта, която показва при кой бюджет всяка позиция влиза в
препоръчания портфейл и при кой излиза от него.

\TODO{среда за измерване: машина, операционна система, компилатор, версия на
пакета за оптимизация; попълва се в деня на измерването}
